\documentclass{article}
\usepackage[margin=1.5cm]{geometry}
\usepackage{amsmath}
\usepackage{enumitem}

\setlist{nosep,leftmargin=*}

\begin{document}
\twocolumn

\title{Chapter 10 In-Class Exercise: Which Lookup Is Faster?}
\author{Prof. Jordan C. Hanson}
\date{}

\maketitle

\section{Build a Large Student Roster}

Suppose student IDs begin at

\begin{verbatim}
first = 2601001
N = 10000
\end{verbatim}
\noindent
Create a list of student names:

\begin{verbatim}
names = ["Student_"+str(i) for i in range(N)]
\end{verbatim}
\noindent
Now represent exactly the same information in two ways.  First, create a \textbf{nested list}:
\begin{verbatim}
student_list = 
[
    [first+i, names[i]] for i in range(N)
]
\end{verbatim}
\noindent
Then create a \textbf{dictionary}:
\begin{verbatim}
student_dict = 
{
    first+i: names[i] for i in range(N)
}
\end{verbatim}
\noindent
Both structures contain the same ID--name pairs.

\section{Lookup Using a Nested List}

To find a student in the nested list, Python must examine records until it finds the correct ID.

\begin{verbatim}
def list_lookup(students,target_id):
    for student in students:
        if student[0] == target_id:
            return student[1]
    return None
\end{verbatim}
\noindent
Test the function:

\begin{verbatim}
target_id = first_id + N - 1
print(list_lookup(student_list,target_id))
\end{verbatim}

\section{Lookup Using a Dictionary}

A dictionary uses the student ID directly as a key:

\begin{verbatim}
def dict_lookup(students,target_id):
    return students.get(target_id)
\end{verbatim}
\noindent
Test it:

\begin{verbatim}
print(dict_lookup(student_dict,target_id))
\end{verbatim}
\noindent
Which method is faster?

\section{Measuring Execution Time}

Python's \texttt{time} module contains a high-resolution timer called \texttt{perf\_counter()}.  Begin with \texttt{import time}.  A single lookup happens so quickly that it may be difficult to measure accurately. Therefore, repeat the lookup many times.

\begin{verbatim}
repeats = 1000
\end{verbatim}
\noindent
Time the nested-list lookup:

\begin{verbatim}
start = time.perf_counter()
for i in range(repeats):
    list_lookup(student_list,target_id)
end = time.perf_counter()
list_time = end - start
\end{verbatim}
\noindent
Now time the dictionary:

\begin{verbatim}
start = time.perf_counter()

for i in range(repeats):
    dict_lookup(student_dict,target_id)
end = time.perf_counter()
dict_time = end - start
\end{verbatim}
\noindent
Display the results:

\begin{verbatim}
print("List:", list_time)
print("Dictionary:", dict_time)
\end{verbatim}

\section{Compare the Results}

Compute how many times faster one method is:

\begin{verbatim}
ratio = list_time / dict_time
print("Ratio:", ratio)
\end{verbatim}
\noindent
Record your results:

\[
t_{\mathrm{list}} = \underline{\hspace{2.2cm}}
\]

\[
t_{\mathrm{dict}} = \underline{\hspace{2.2cm}}
\]

\[
\frac{t_{\mathrm{list}}}
     {t_{\mathrm{dict}}}
=
\underline{\hspace{2.2cm}}
\]
\noindent
Repeat the experiment using

\begin{verbatim}
N = 100
N = 1000
N = 10000
\end{verbatim}

For each value of \(N\), search for the \textbf{last} student.

\end{document}
